\documentclass[11pt, letterpaper]{article}

\usepackage[margin=2cm]{geometry}
\usepackage[spanish]{babel}
\usepackage{amsmath,amssymb}
\usepackage{xcolor}
\usepackage{fancyhdr}
\usepackage{tcolorbox}
\usepackage{enumitem}

\definecolor{primary}{RGB}{255, 107, 53}
\definecolor{darkslate}{RGB}{15, 23, 42}

\pagestyle{fancy}
\fancyhf{}
\fancyhead[L]{\small \textbf{Universidad Nacional} -- Departamento de Matemáticas}
\fancyhead[R]{\small \textbf{Cálculo III} -- Semestre 2026-1}
\fancyfoot[C]{\thepage}

\begin{document}

% MEMBRETE Y TABLA DE IDENTIFICACIÓN
\begin{tcolorbox}[colback=darkslate, colframe=primary, arc=2mm, boxrule=1.5pt]
  \color{white}
  \begin{center}
    {\Large \bfseries EVALUACIÓN PARCIAL N$^\circ$ 2 -- CÁLCULO VECTORIAL}\\[0.2cm]
    {\small Tiempo máximo: 90 minutos \quad $\bullet$ \quad Puntaje total: 60 puntos \quad $\bullet$ \quad Exigencia: 60\%}
  \end{center}
\end{tcolorbox}

\vspace{0.3cm}

\noindent
\begin{tabular}{|p{10cm}|p{6.8cm}|}
  \hline
  \textbf{Nombre del Estudiante:} & \textbf{RUT / Identificación:} \\
  \vspace{0.3cm} & \vspace{0.3cm} \\
  \hline
  \textbf{Profesor:} Dr. M. Fernández & \textbf{Fila / Sección:} 101 \\
  \hline
\end{tabular}

\vspace{0.4cm}

\begin{tcolorbox}[colback=primary!10, colframe=primary, arc=2mm, title=\textbf{Instrucciones Generales}]
  $\bullet$ Responda con letra clara y orden en el espacio indicado. Justifique cada paso matemáticamente. \\
  $\bullet$ Se prohíbe el uso de dispositivos celulares o calculadoras programables.
\end{tcolorbox}

\vspace{0.5cm}

\subsection*{Problema 1 (15 Puntos) -- Integrales de Línea y Campos Conservativos}
Considere el campo vectorial bidimensional:
\begin{equation*}
  \mathbf{F}(x, y) = \left( 2x y^3 + e^x \sin y \right) \mathbf{i} + \left( 3x^2 y^2 + e^x \cos y \right) \mathbf{j}
\end{equation*}

\begin{enumerate}[label=\textbf{\alph*)}]
  \item Demuestre que el campo $\mathbf{F}$ es conservativo en todo $\mathbb{R}^2$.
  \item Encuentre la función potencial escalar $\phi(x, y)$ tal que $\mathbf{F} = \nabla \phi$.
  \item Calcule el trabajo realizado por $\mathbf{F}$ al desplazar una partícula a lo largo de la curva suave $C$ desde el punto $A(0, 0)$ hasta $B(1, \pi/2)$.
\end{enumerate}

\vspace{3.5cm}

\subsection*{Problema 2 (20 Puntos) -- Teorema de Green en el Plano}
Utilice el Teorema de Green para evaluar la integral de línea en sentido antihorario:
\begin{equation*}
  \oint_C \left( y^2 - \arctan x \right) dx + \left( 3x - \ln(y^2 + 1) \right) dy
\end{equation*}
Donde $C$ es la frontera de la región encerrada entre las curvas $y = x^2$ e $y = 4$.

\vspace{3.5cm}

\subsection*{Problema 3 (25 Puntos) -- Teorema de la Divergencia de Gauss}
Sea $\mathbf{F}(x, y, z) = x^3 \mathbf{i} + y^3 \mathbf{j} + z^3 \mathbf{k}$ y sea $S$ la superficie de la esfera unitaria $x^2 + y^2 + z^2 = 1$ orientada con la normal exterior. Calcule el flujo total:
\begin{equation*}
  \Phi = \iint_S \mathbf{F} \cdot d\mathbf{S}
\end{equation*}

\end{document}
